\section{Discussion}

The theorem separates representation from performance. It removes arbitrary-support frames from exact optimization without asserting that a direct enumeration is practically efficient. The adverse runtime study illustrates why this distinction matters: a polynomial exponent and a smaller normal form do not automatically outperform a factorized exact referee on finite workloads.

The sharpness witness also identifies the obstruction to support one. At support two, no nonempty proper subset need preserve both frame anticommutation and the shared-label syndrome. A locally removable frame letter can therefore be globally coupled to the shared operator. The upper proof and lower witness meet at this parity boundary rather than at an empirical cutoff.

The result is useful primarily as a formalization of one cost model. It supplies a complete representational reduction, an exact algorithmic upper bound, and a reproducible negative performance result. It does not establish a new universal synthesis primitive. Practical significance would require either a better algorithm exploiting the normal form, a broader theorem, or an independently evaluated resource improvement under a physically meaningful model.
